\documentclass[conference]{IEEEtran}
\IEEEoverridecommandlockouts
% The preceding line is only needed to identify funding in the first footnote. If that is unneeded, please comment it out.
\usepackage{cite}
\usepackage{amsmath,amssymb,amsfonts}
\usepackage{algorithmic}
\usepackage{graphicx}
\usepackage{textcomp}
\usepackage{xcolor}
\def\BibTeX{{\rm B\kern-.05em{\sc i\kern-.025em b}\kern-.08em
    T\kern-.1667em\lower.7ex\hbox{E}\kern-.125emX}}

\begin{document}

\title{Hero-Dash: A Multimodal Serious Game for Inclusive Disaster Early Warning Training in Hearing-Impaired Children\\
}

\author{
\IEEEauthorblockN{T. L. Samarasekara\IEEEauthorrefmark{1}, H. H. Liyanage\IEEEauthorrefmark{1}, W. D. T. N. Weerasinghe\IEEEauthorrefmark{1}, \\ W. A. T. R. Wijesekara\IEEEauthorrefmark{1}, Samanthi Rubasin Siriwardana\IEEEauthorrefmark{2}, Lokesha Weerasinghe\IEEEauthorrefmark{1}}
\IEEEauthorblockA{\IEEEauthorrefmark{1}Sri Lanka Institute of Information Technology, Colombo, Sri Lanka\\
Email: \{it22337726, it22550262, it22586148, it22322630\}@my.sliit.lk, lokesha.w@sliit.lk}
\IEEEauthorblockA{\IEEEauthorrefmark{2}York St John University, York, United Kingdom\\
Email: s.rubasinsiriwarda@yorksj.ac.uk}
}

\maketitle

\begin{abstract}
Hearing-impaired children are particularly vulnerable during emergencies due to the inaccessibility of traditional early warning systems. To address this critical gap, this paper presents Hero-Dash, an adaptive, web-based serious game designed to train children (ages 5--14) to identify and respond to critical disaster warnings. The system integrates a 3D immersive simulation with a multi-algorithm machine learning pipeline to personalize the training experience. Through narrative-driven missions, children learn to recognize five key warning sounds: tsunami, earthquake, flood, air raid, and fire alarms. The backend utilizes Bayesian Knowledge Tracing to monitor skill mastery, Item Response Theory to assess crisis readiness, and Thompson Sampling to optimize scenario selection. Furthermore, to accommodate varying degrees of hearing loss, the system incorporates personalized audio equalization based on individual audiograms, alongside visual and vibrotactile alerts. Initial development and usability testing demonstrate high user engagement and system feasibility. By providing a scalable, disability-inclusive training platform, Hero-Dash offers a practical technological approach to fulfilling the Sendai Framework's mandate for inclusive disaster risk reduction.
\end{abstract}

\begin{IEEEkeywords}
disaster early warning, hearing impairment, serious games, adaptive machine learning, accessibility, Bayesian Knowledge Tracing
\end{IEEEkeywords}

\section{Introduction}
The Sendai Framework for Disaster Risk Reduction \cite{b1} emphasizes the critical need for disability-inclusive early warning systems. Despite this global mandate, approximately 34 million hearing-impaired children worldwide remain disproportionately vulnerable during crises due to a lack of accessible alert infrastructure and tailored preparedness training \cite{b2}. While serious games have proven effective in general emergency preparedness---yielding significantly faster evacuation decision times \cite{b3}---and adaptive machine learning (ML) has consistently improved educational outcomes \cite{b4}, these technologies have not been successfully combined to address the specific needs of hearing-impaired learners \cite{b5}.

Hearing-impaired children face compounded barriers in emergency situations. Central Auditory Processing deficits often impair their ability to discriminate warning sirens from background noise, even when utilizing hearing aids \cite{b6}. Furthermore, current digital training tools lack the necessary AI-driven personalization to adapt to individual audiological profiles and learning paces \cite{b7}. There remains a stark disconnect between the proven efficacy of adaptive serious games and the complete lack of deployment of such tools for disability-inclusive crisis preparedness.

This paper introduces Hero-Dash, a novel serious game designed specifically for disaster early warning training in hearing-impaired children. By integrating narrative-driven gameplay, a six-algorithm adaptive ML pipeline, and rigorous psychometric measurement, Hero-Dash aims to bridge the gap between inclusive policy mandates and practical technological solutions.

\section{System Architecture and Design}

\subsection{Interactive Environment and Gamification}
Hero-Dash is deployed as a client-server web application. The frontend utilizes ReactJS and Three.js to render a procedurally generated 3D urban environment, managed by Zustand state management. The backend, built with FastAPI and SQLAlchemy, manages the ML algorithms and psychometric engine. 

To maintain engagement, stories frame crisis responses as rescue missions (e.g., ``Tsunami Evacuation: Help families reach high ground''). This bilingual narrative (English and Sinhala) engages children through goal-setting and achievement unlocks, operationalizing goal-setting theory \cite{b8} and narrative transportation \cite{b9}. During gameplay, crisis scenarios are triggered at dynamically determined intervals. Each scenario requires a specific life-saving response: moving to high ground (right) for a tsunami, executing a full stop for an earthquake, seeking designated safe zones for floods, staying centered for air raids, and evacuating left for fire alarms.

\subsection{Multimodal Alert Delivery}
To ensure accessibility across the spectrum of hearing impairment, alerts are transmitted simultaneously across three sensory channels:

\begin{enumerate}
    \item \textbf{Audio Channel:} Sound profiles are customized using a 6-band parametric equalizer (250--8000 Hz) applying NAL-NL2 gains computed directly from the user's individual audiogram. Dynamic range compression prevents discomfort, while clinician-configurable environmental noise (0.0--1.0) trains figure-ground separation.
    \item \textbf{Visual Channel:} The interface employs WCAG 2.1 AAA compliant color-coded overlays, animated directional arrows, a countdown timer, and bilingual text labels. Visual prompts dynamically fade as the system detects skill mastery.
    \item \textbf{Haptic Channel:} For users with profound hearing loss, the system leverages the Web Vibration API to trigger distinct, crisis-specific vibrotactile patterns (e.g., a tsunami alert uses a rapid quadruple pulse).
\end{enumerate}

\section{Adaptive Machine Learning Pipeline}

Rather than relying on a single monolithic model or black-box deep learning, Hero-Dash employs a priority-ordered cascade of six distinct algorithms. This modular, interpretable approach ensures clinical transparency and allows the system to address specific pedagogical requirements in real-time.

\begin{enumerate}
    \item \textbf{Scenario Selection:} Thompson Sampling manages the exploration-exploitation tradeoff, selecting scenarios based on historical learning gains.
    \item \textbf{Mastery Tracking:} Bayesian Knowledge Tracing (BKT) uses Hidden Markov Models to monitor the acquisition of five critical auditory skills: frequency discrimination, temporal pattern recognition, figure-ground separation, sound-action mapping, and sustained auditory attention. Mastery is declared when the probability of learning passes 0.95.
    \item \textbf{Readiness Estimation:} A 2-Parameter Logistic Item Response Theory (IRT) model dynamically estimates the user's overall crisis readiness, adapting item difficulty to the user's latent ability.
    \item \textbf{Retention Management:} An SM-2 spaced repetition algorithm schedules the review of previously mastered scenarios to prevent memory decay over time.
    \item \textbf{Cognitive Load and Flow Regulation:} Real-time monitors track cognitive load (via reaction time variance and error clusters) and flow-state (maintaining success rates between 60--85\%). These adjust environmental difficulty parameters, such as background noise, without altering the core scenario type.
\end{enumerate}

The algorithms operate on a 10-second polling cycle. The cascade prioritizes BKT remediation and spaced repetition reviews before utilizing Thompson Sampling for exploration, ensuring that immediate pedagogical needs are met before introducing new challenges.

\section{Psychometric Validation Framework}

To ensure the training provides measurable clinical value, Hero-Dash incorporates a comprehensive psychometric evaluation suite aligned with American Educational Research Association standards \cite{b10}. The system occasionally transitions from the gamified ``Mission Mode'' to a standardized ``Assessment Mode''---a fixed 20-trial protocol with adaptive ML disabled---to capture unassisted performance data.

The framework automatically calculates internal consistency (Cronbach's $\alpha \geq 0.80$) and Test-Retest reliability between consecutive standardized sessions. It computes the Standard Error of Measurement (SEM) and the Minimal Detectable Change (MDC), allowing clinicians to confidently distinguish genuine learning progress from standard measurement variance \cite{b12}. Furthermore, assessment items are conceptually mapped to established clinical instruments, including SCAN-3:CH (auditory processing) \cite{b13} and CHAPPS (functional listening) \cite{b14}, supporting construct validity. 

\section{Implementation and Results}

\subsection{Development Outcomes and Usability}
The end-to-end system architecture has been fully implemented. This includes the 3D frontend, the algorithmic backend, and a comprehensive therapist dashboard providing skill trajectory visualization and pre/post comparison metrics. The system successfully processes algorithmic cascades within a 312 ms latency average, well under the 500 ms budget required for seamless real-time adaptation.

Initial usability testing was conducted with 23 hearing-impaired children (ages 5--14). The results demonstrated a mean session engagement rate of 87\%, which is a substantial improvement over the 60\% engagement typical of traditional clinical auditory training exercises \cite{b15}.

\subsection{Machine Learning Performance Observations}
During development testing (comprising 47 sessions with proxy users to evaluate system stability), the BKT module successfully modeled skill acquisition trajectories. Notably, figure-ground separation consistently showed the slowest acquisition rate, reaching a mean mastery probability of only 0.55 at 30 trials, compared to 0.87 for frequency discrimination. This aligns with existing clinical literature regarding the specific difficulties hearing-impaired learners face with background noise \cite{b11}.

Additionally, the IRT models showed a positive monotonic growth in crisis readiness across multiple sessions. Mean $\theta$ estimates increased from -0.42 (below-average readiness) at session 1 to +1.31 (above-average readiness) by session 6, validating the adaptive testing progression. The SM-2 spaced repetition module maintained a review interval adherence rate of 91.5\%.

\begin{figure}[htbp]
\centering
\includegraphics[width=0.8\columnwidth]{figures/gameplay.png}
\caption{Hero-Dash 3D driving environment with active tsunami crisis alert. The multimodal overlay is displayed simultaneously with the siren audio and haptic vibration pattern. The HUD shows real-time BKT mastery estimates and IRT readiness score.}
\label{fig:gameplay}
\end{figure}

\begin{table}[htbp]
\centering
\caption{Hero-Dash therapist dashboard: Child P-007 Session Summary}
\label{tab:dashboard}
\small
\begin{tabular}{|l|l|}
\hline
\textbf{Patient Information} & \\
\hline
Child ID & P-007 \\
Age & 9 years \\
Hearing Loss & Moderate \\
Session Date & 2026-02-15 \\
Session Duration & 14:32 \\
\hline
\textbf{Skill Mastery Trajectories} & \textbf{BKT P(L)} \\
\hline
Frequency Discrimination & 0.87\,\checkmark \\
Temporal Pattern Recognition & 0.71 \\
Figure-Ground Separation & 0.55\,! \\
Sound-Action Mapping & 0.91\,\checkmark \\
Auditory Attention Duration & 0.66 \\
\hline
\textbf{IRT Crisis Readiness} & $\theta = 1.24$ (SE = 0.31) \\
\hline
\textbf{Psychometric Summary (Session 2)} & \\
\hline
Cronbach's $\alpha$ & 0.83 \\
Standard Error of Measurement (SEM) & 4.2 \\
Minimal Detectable Change$_{95}$ (MDC$_{95}$) & 11.6 \\
Pre-Assessment Accuracy & 48\% \\
Post-Assessment Accuracy & 71\% \\
Cohen's $d$ Effect Size & 0.74 \\
\hline
\textbf{SM-2 Spaced Repetition} & \\
\hline
Next Review Item & Figure-Ground Separation \\
Review Due & +3 days \\
\hline
\end{tabular}
\end{table}

\begin{figure}[htbp]
\centering
\includegraphics[width=0.8\columnwidth]{figures/dashboard.png}
\caption{Therapist dashboard screenshot showing BKT skill mastery trajectories, IRT crisis readiness score with standard error, and session-level psychometric summary statistics.}
\label{fig:dashboard}
\end{figure}

\section{Discussion}

\subsection{Operationalizing Inclusive Frameworks}
Hero-Dash successfully translates the abstract goals of the Sendai Framework \cite{b1} into a deployable technological solution. By moving beyond traditional audio-only sirens to provide multimodal, audiogram-personalized alerts, the system offers a scalable method for improving disaster resilience among a highly vulnerable population. The combination of serious game design with clinical-grade psychometrics ensures that the tool is both engaging for children and valid for audiologists and therapists.

\subsection{Limitations}
While development-phase testing confirms system stability and high user engagement, a primary limitation is the current lack of empirical efficacy data from a large-scale clinical trial with the target demographic. Furthermore, the BKT and IRT parameters currently rely on literature-derived priors and expert judgment; recalibration using empirical data from the target population will be necessary. Finally, the Web Audio API provides slightly less precise audiological control than dedicated clinical hearing aid fitting software, which may limit equalization accuracy for users with extreme audiogram configurations (e.g., precipitous high-frequency loss).

\subsection{Future Work}
The immediate next step is a pre-registered Randomized Controlled Trial (RCT) involving 60 hearing-impaired children across Sri Lankan clinical sites to rigorously evaluate learning outcomes against waitlist controls. Primary outcomes will include SCAN-3:CH composite scores and emergency response accuracy. Future technical iterations will aim to expand the crisis sound library to include severe weather and industrial hazards, integrate the platform with emerging smart-city early warning infrastructures, and pursue ISO 13485 medical device regulatory pathways.

\section{Conclusion}

This paper presents Hero-Dash, an innovative serious game that leverages multimodal sensory delivery and an interpretable machine learning pipeline to provide accessible disaster preparedness training for hearing-impaired children. By combining narrative engagement with rigorous, clinically-aligned psychometric assessment, the platform demonstrates the technical feasibility of personalized, disability-inclusive early warning education. Development-phase implementation confirms algorithmic stability, high engagement, and system usability. Following planned clinical trials to establish empirical efficacy, Hero-Dash provides a concrete pathway toward operationalizing the Sendai Framework, with the potential to significantly enhance the safety and resilience of millions of vulnerable children globally.

\begin{thebibliography}{00}
\bibitem{b1} UNDRR, \textit{Sendai Framework for Disaster Risk Reduction 2015-2030}, United Nations Office for Disaster Risk Reduction, Geneva, 2015.
\bibitem{b2} World Health Organization, \textit{World Report on Hearing}, WHO, Geneva, 2021.
\bibitem{b3} R. Lovreglio, E. Ronchi, and D. Nilsson, ``A model of the decision-making process during pre-evacuation,'' \textit{Fire Safety Journal}, vol. 78, pp. 168-179, 2021.
\bibitem{b4} L. Chittaro and F. Buttussi, ``Serious games for emergency preparedness: Evaluation of an interactive vs. a non-interactive simulation of a terror attack,'' \textit{Computers in Human Behavior}, vol. 139, art. 107535, 2023.
\bibitem{b5} X. Wang, R. T. Huang, M. Sommer, and B. Pei, ``The efficacy of artificial intelligence-enabled adaptive learning systems from 2010 to 2022: A meta-analysis,'' \textit{Journal of Educational Computing Research}, vol. 62, no. 4, 2024.
\bibitem{b6} F. E. Musiek, G. D. Chermak, and J. Weihing, \textit{Handbook of Central Auditory Processing Disorder}, 2nd ed., Plural Publishing, San Diego, 2020.
\bibitem{b7} J. Zhao, I. Neto, A. C. Pires, and C. Tome-Pires, ``Digital technologies for deaf and hard of hearing children: A systematic review,'' in \textit{Proc. ACM CHI 2025}, ACM, New York, 2025.
\bibitem{b8} E. A. Locke and G. P. Latham, \textit{New Developments in Goal Setting and Task Performance}. New York: Routledge, 2013.
\bibitem{b9} M. C. Green, ``Narrative and storytelling in entertainment and pedagogy,'' in \textit{The Oxford Handbook of Media Psychology}, D. Zillmann, J. Bryant, and A. C. Huston, Eds. New York: Oxford University Press, 2008, pp. 237--256.
\bibitem{b10} AERA, ``Standards for Educational and Psychological Testing,'' American Educational Research Association, Washington, DC, 2014.
\bibitem{b11} Q. Liu, Y. Zhuang, H. Bi, Z. Huang, W. Huang, and J. Li, ``Survey of computerized adaptive testing: A machine learning perspective,'' \textit{IEEE Trans. Learn. Technol.}, vol. 17, pp. 1--18, 2024.
\bibitem{b12} C. Terwee, et al., ``Quality criteria were proposed for measurement properties of health status questionnaires,'' \textit{J. Clin. Epidemiol.}, vol. 60, no. 1, pp. 34--42, 2007.
\bibitem{b13} R. W. Keith, \textit{SCAN-3:C Tests for Auditory Processing Disorders in Children}, Pearson, San Antonio, TX, 2009.
\bibitem{b14} W. Smoski, M. Brunt, and J. Tannahill, \textit{Children's Auditory Performance Scale (CHAPPS)}, Educational Audiology Association, 1998.
\bibitem{b15} L. M. Stough and D. Kang, ``The Sendai Framework for Disaster Risk Reduction and persons with disabilities,'' \textit{International Journal of Disaster Risk Science}, vol. 6, no. 2, pp. 140-149, 2015.
\end{thebibliography}

\end{document}